\documentclass[UTF8]{ctexart}
\usepackage{xeCJK}
\usepackage{geometry}
\geometry{left=2cm,right=2cm,top=2.5cm,bottom=2.5cm}
\setlength{\headheight}{12.7pt}

\usepackage{titlesec}
\titleformat{\section}
  {\normalfont\large\bfseries}
  {\thesection}
  {1em}
  {}
\titlespacing*{\section}
  {0pt}
  {3.5ex plus 1ex minus .2ex}
  {2.3ex plus .2ex}

\usepackage{amsmath}
\usepackage{amssymb}
\usepackage{amsthm}
\usepackage{enumitem}
\setlist[enumerate]{itemsep=0pt,topsep=0pt,parsep=0pt,leftmargin=0pt,itemindent=2.5em}
\setlist[itemize]{itemsep=0pt,topsep=0pt,parsep=0pt,leftmargin=0pt,itemindent=2.5em}
\usepackage{fancyhdr}
\pagestyle{fancy}
\fancyhf{}
\usepackage{graphicx}
\usepackage{hyperref}
\usepackage{indentfirst}
\usepackage{lmodern}


\begin{document}
\setlength{\abovedisplayskip}{1pt}
\setlength{\belowdisplayskip}{1pt}

\section{while 程序指称语义}

\hypertarget{sec:定义1.1}{}
\noindent\textbf{定义1.1 (指称语义)}

程序$S$的\textbf{部分正确性语义}定义为
\[
    \mathcal{M}[\![S]\!]:\Sigma\to\mathcal{P}(\Sigma),\quad\sigma\mapsto
    \{\,\tau\,|\,\langle S,\sigma\rangle\text{终止于}\tau\,\},
\]
其\textbf{完全正确性语义}定义为
\[
    \mathcal{M}_{\mathrm{tot}}[\![S]\!]:\Sigma\to\mathcal{P}(\Sigma\cup\{\bot\}),
    \quad\sigma\mapsto\mathcal{M}[\![S]\!](\sigma)\cup
    \{\,\bot\,|\,\langle S,\sigma\rangle\text{发散}\,\}.
\]

\vspace{10pt}

当$\langle S,\sigma\rangle$终止于$\tau$时, $\mathcal{M}[\![S]\!](\sigma)=
\mathcal{M}_{\mathrm{tot}}[\![S]\!](\sigma)=\{\tau\}$,

当$\langle S,\sigma\rangle$发散时,
$\mathcal{M}[\![S]\!](\sigma)=\varnothing$,
$\mathcal{M}_{\mathrm{tot}}[\![S]\!](\sigma)=\{\bot\}$.

\noindent{\kaishu 备注.}

\begin{enumerate}
    \item 把$\bot$视为一个始终发散的扩展状态,
    从而把上述两个语义映射的定义域扩充到$\Sigma\cup\{\bot\}$.
    \vspace{3pt}
    \item 对$X\subset\Sigma\cup\{\bot\}$, 定义
    $\displaystyle\mathcal{N}[\![S]\!](X)=
    \bigcup_{\sigma\in X}\mathcal{N}[\![S]\!](\sigma)$,
    这里$\mathcal{N}$指代$\mathcal{M}$或$\mathcal{M}_{\mathrm{tot}}$
    (下文同理). 显然这两个函数是单调的.
\end{enumerate}

\vspace{10pt}

\hypertarget{sec:定理1.1}{}
\noindent\textbf{定理1.1 (seq语句的指称语义)}

\begin{enumerate}
    \item 对任意的程序$S_1,S_2$和$X\subset\Sigma\cup\{\bot\}$, 有
    $\mathcal{N}[\![S_1;S_2]\!](X)=\mathcal{N}[\![S_2]\!]\,
    \mathcal{N}[\![S_1]\!](X)$,
    \item 对任意的程序$S_1,S_2,S_3$和$X\subset\Sigma\cup\{\bot\}$, 有
    $\mathcal{N}[\![(S_1;S_2);S_3]\!](X)=
    \mathcal{N}[\![S_1;(S_2;S_3)]\!](X)$.
\end{enumerate}

\noindent{\kaishu 证明.}

\begin{enumerate}
    \item 任取$\sigma\in\Sigma\cup\{\bot\}$和
    从$\langle S_1,\sigma\rangle$开始的计算$\langle S_1\,\!\!^i,\sigma^i\rangle$,
    归纳易证$\langle S_1\,\!\!^i;S_2,\sigma^i\rangle$是一个转移序列.

    \hspace{2.17em}
    如果$\langle S_1,\sigma\rangle$终止于$\sigma'$
    且$\langle S_2,\sigma'\rangle$终止于$\tau$, 那么
    \[
        \langle S_1,\sigma\rangle\to^*\langle\mathrm{E},\sigma'\rangle
        \quad\implies\quad
        \langle S_1;S_2,\sigma\rangle\to^*\langle\mathrm{E};S_2,\sigma'\rangle=
        \langle S_2,\sigma'\rangle\to^*\langle\mathrm{E},\tau\rangle,
    \]
    此时
    \[\left\{\begin{aligned}
        &\mathcal{M}[\![S_1;S_2]\!](\sigma)=\{\tau\}=
        \mathcal{M}[\![S_2]\!](\sigma')=
        \mathcal{M}[\![S_2]\!]\mathcal{M}[\![S_1]\!](\sigma),\\
        &\mathcal{M}_{\mathrm{tot}}[\![S_1;S_2]\!](\sigma)=\{\tau\}=
        \mathcal{M}_{\mathrm{tot}}[\![S_2]\!](\sigma')=
        \mathcal{M_{\mathrm{tot}}}[\![S_2]\!]
        \mathcal{M}_{\mathrm{tot}}[\![S_1]\!](\sigma),\\
    \end{aligned}\right.\]

    \hspace{2.17em}
    如果$\langle S_1,\sigma\rangle$终止于$\sigma'$
    且$\langle S_2,\sigma'\rangle$发散, 那么
    \[
        \langle S_1,\sigma\rangle\to^*\langle\mathrm{E},\sigma'\rangle
        \quad\implies\quad
        \langle S_1;S_2,\sigma\rangle\to^*\langle\mathrm{E};S_2,\sigma'\rangle=
        \langle S_2,\sigma'\rangle\text{发散},
    \]
    此时
    \[\left\{\begin{aligned}
        &\mathcal{M}[\![S_1;S_2]\!](\sigma)=\varnothing=
        \mathcal{M}[\![S_2]\!](\sigma')=
        \mathcal{M}[\![S_2]\!]\mathcal{M}[\![S_1]\!](\sigma),\\
        &\mathcal{M}_{\mathrm{tot}}[\![S_1;S_2]\!](\sigma)=\{\bot\}=
        \mathcal{M}_{\mathrm{tot}}[\![S_2]\!](\sigma')=
        \mathcal{M_{\mathrm{tot}}}[\![S_2]\!]
        \mathcal{M}_{\mathrm{tot}}[\![S_1]\!](\sigma),\\
    \end{aligned}\right.\]

    \hspace{2.17em}
    如果$\langle S_1,\sigma\rangle$发散(包括$\sigma=\bot$), 那么
    $\langle S_1;S_2,\sigma\rangle$也发散, 此时
    \vspace{3pt}
    \[\left\{\begin{aligned}
        &\mathcal{M}[\![S_1;S_2]\!](\sigma)=\varnothing=
        \mathcal{M}[\![S_2]\!](\varnothing)=
        \mathcal{M}[\![S_2]\!]\mathcal{M}[\![S_1]\!](\sigma),\\
        &\mathcal{M}_{\mathrm{tot}}[\![S_1;S_2]\!](\sigma)=\{\bot\}=
        \mathcal{M}_{\mathrm{tot}}[\![S_2]\!](\bot)=
        \mathcal{M_{\mathrm{tot}}}[\![S_2]\!]
        \mathcal{M}_{\mathrm{tot}}[\![S_1]\!](\sigma),\\
    \end{aligned}\right.\]

    \vspace{3pt}
    \item \hspace*{\fill}$\begin{aligned}[t]
        \mathcal{N}[\![(S_1;S_2);S_3]\!](X)=\,&
        \mathcal{N}[\![S_3]\!]\mathcal{N}[\![S_1;S_2]\!](X)=
        \mathcal{N}[\![S_3]\!]\mathcal{N}[\![S_2]\!]\mathcal{N}[\![S_1]\!](X)\\
        =\,&\mathcal{N}[\![S_2;S_3]\!]\mathcal{N}[\![S_1]\!](X)=
        \mathcal{N}[\![S_1;(S_2;S_3)]\!](X).
    \end{aligned}$ \hspace*{\fill}\null
\end{enumerate}





\newpage

\hypertarget{sec:定理1.2}{}
\noindent\textbf{定理1.2 (if语句的指称语义)}

对任意的$B:\mathrm{Bool}$, 程序$S_1,S_2$和$X\subset\Sigma\cup\{\bot\}$, 有
$\mathcal{N}[\![\,\mathbf{if}\,\,B\,\,\mathbf{then}\,\,S_1\,\,
\mathbf{else}\,\,S_2\,\,\mathbf{fi}\,]\!](X)=$
\[
    \mathcal{N}[\![S_1]\!](X\cap[\![B]\!])\,\cup\,
    \mathcal{N}[\![S_2]\!](X\cap[\![\neg\,B]\!])\,\cup\,
    \{\,\bot\,\mid\,\bot\in X\land\mathcal{N}=\mathcal{M}_{\mathrm{tot}}\,\}.
\]

\noindent{\kaishu 证明.}

任取$\sigma\in\Sigma\cup\{\bot\}$.
如果$\sigma\models B$, 那么$\langle
\mathbf{if}\,\,B\,\,\mathbf{then}\,\,S_1\,\,\mathbf{else}\,\,S_2\,\,\mathbf{fi},
\sigma\rangle\to\langle S_1,\sigma\rangle$, 所以
\[
    \mathcal{N}[\![\,\mathbf{if}\,\,B\,\,\mathbf{then}\,\,S_1\,\,
    \mathbf{else}\,\,S_2\,\,\mathbf{fi}\,]\!](\sigma)=
    \mathcal{N}[\![S_1]\!](\sigma),
\]
如果$\sigma\not\models B$, 那么$\langle
\mathbf{if}\,\,B\,\,\mathbf{then}\,\,S_1\,\,\mathbf{else}\,\,S_2\,\,\mathbf{fi},
\sigma\rangle\to\langle S_2,\sigma\rangle$, 所以
\[
    \mathcal{N}[\![\,\mathbf{if}\,\,B\,\,\mathbf{then}\,\,S_1\,\,
    \mathbf{else}\,\,S_2\,\,\mathbf{fi}\,]\!](\sigma)=
    \mathcal{N}[\![S_2]\!](\sigma),
\]
如果$\sigma=\bot$, 那么
\begin{equation*}
    \mathcal{N}[\![\,\mathbf{if}\,\,B\,\,\mathbf{then}\,\,S_1\,\,
    \mathbf{else}\,\,S_2\,\,\mathbf{fi}\,]\!](\sigma)=
    \{\,\bot\,\mid\,\mathcal{N}=\mathcal{M}_{\mathrm{tot}}\}.
\tag*{\qedsymbol}\end{equation*}

\vspace{10pt}

\hypertarget{sec:定理1.3}{}
\noindent\textbf{定理1.3 (while语句的指称语义)}

对任意的$B:\mathrm{Bool}$和程序$S$, 递归定义程序
\[A^k=\left\{\begin{aligned}
    &\mathbf{while}\,\,\mathrm{true}\,\,\mathbf{do}\,\,\mathrm{skip}
    \,\,\mathbf{od},\quad k=0,\\
    &\mathbf{if}\,\,B\,\,\mathbf{then}\,\,S;A^{k-1}\,\,\mathbf{else}
    \,\,\mathrm{skip}\,\,\mathbf{fi},\quad k>0,
\end{aligned}\right.\]
则对任意的$X\subset\Sigma\cup\{\bot\}$有
\[
    \mathcal{M}[\![\,\mathbf{while}\,\,B\,\,\mathbf{do}
    \,\,S\,\,\mathbf{od}\,]\!](X)=\bigcup_k\mathcal{M}[\![A^k]\!](X).
\]


\noindent{\kaishu 证明.}

记$W=\mathbf{while}\,\,B\,\,\mathbf{do}\,\,S\,\,\mathbf{od}$.
第一, 对$k$归纳证明, 对任意$\sigma\in\Sigma\cup\{\bot\}$有
$\mathcal{M}[\![A^k]\!](\sigma)\subset\mathcal{M}[\![W]\!](\sigma)$.

显然$k=0$时$\mathcal{M}[\![A^0]\!](\sigma)=\varnothing$\,;
假设命题对$k$成立, 任取$\sigma$, 那么
\begin{enumerate}
    \item 如果$\sigma\models B$, 此时$\langle A^{k+1},\sigma\rangle\to
    \langle S;A^k,\sigma\rangle$终止, 从而$\langle S,\sigma\rangle$
    一定终止, 设其终止于$\sigma'$, 则
    \[
        \langle A^{k+1},\sigma\rangle\to\langle S;A^k,\sigma\rangle\to^*
        \langle A^k,\sigma'\rangle,\quad
        \langle W,\sigma\rangle\to\langle S;W,\sigma\rangle\to^*
        \langle W,\sigma'\rangle,
    \]
    从而$\mathcal{M}[\![A^{k+1}]\!](\sigma)=\mathcal{M}[\![A^k]\!](\sigma')\subset
    \mathcal{M}[\![W]\!](\sigma')=\mathcal{M}[\![W]\!](\sigma)$.
    \item 如果$\sigma\not\models B$, 则$\mathcal{M}[\![A^{k+1}]\!](\sigma)=
    \{\sigma\}=\mathcal{M}[\![W]\!](\sigma)$.
    \item 如果$\sigma=\bot$, 则$\mathcal{M}[\![A^{k+1}]\!](\sigma)=\varnothing$.
\end{enumerate}

综上归纳成立, 对任意的$k$和$\sigma$有
$\mathcal{M}[\![A^k]\!](\sigma)\subset\mathcal{M}[\![W]\!](\sigma)$,
从而$\displaystyle\bigcup_k\mathcal{M}[\![A^k]\!](\sigma)\subset
\mathcal{M}[\![W]\!](\sigma)$.

\vspace{3pt}
第二, 对任意的$\sigma$和$\tau\in\mathcal{M}[\![W]\!](\sigma)$,
记该计算中$W$出现次数为$n\geqslant 1$,
对$n$归纳证明$\tau\in\mathcal{M}[\![A^n]\!](\sigma)$.

首先, 如果计算中$W$只出现$1$次, 则一定有$\sigma\not\models B$, 即
$\tau=\sigma\in\mathcal{M}[\![A^1]\!](\sigma)$;

其次假设命题对$n$成立, 并且计算中$W$出现了$n+1$次, 则一定有$\sigma\models B$.
此时$\langle W,\sigma\rangle\to\langle S;W,\sigma\rangle$终止,
从而$\langle S,\sigma\rangle$一定终止, 设其终止于$\sigma'$, 则
\[
    \langle W,\sigma\rangle\to\langle S;W,\sigma\rangle\to^*
    \langle W,\sigma'\rangle,\quad
    \langle A^{n+1},\sigma\rangle\to\langle S;A^n,\sigma\rangle\to^*
    \langle A^n,\sigma'\rangle
\]
易证$\langle W,\sigma'\rangle$的计算中$W$出现$n$次, 依假设有
$\tau\in\mathcal{M}[\![A^n]\!](\sigma')=\mathcal{M}[\![A^{n+1}]\!](\sigma)$.

\vspace{3pt}
综上归纳成立, 对任意的$\sigma$有$\displaystyle
\mathcal{M}[\![W]\!](\sigma)\subset\bigcup_k\mathcal{M}[\![A^k]\!](\sigma)$.
\hfill $\qedsymbol$

\vspace{10pt}

\hypertarget{sec:定理1.4}{}
\noindent\textbf{定理1.4}

\begin{enumerate}
    \item 任给$\tau\in\mathcal{M}[\![S]\!](\sigma)$,
    则$\sigma,\tau$在$\mathrm{Var}-\mathrm{change}(S)$上相同.
    \item 任给$\tau\in\mathcal{M}[\![S]\!](\sigma)$
    和$\tau'\in\mathcal{M}[\![S]\!](\sigma')$,
    如果$\sigma,\sigma'$在$\mathrm{var}(S)$上相同, 则$\tau,\tau'$也如此.
\end{enumerate}

\end{document}